%
% 000-session9.tex
% Presentation Session 9: Practice 2 - Client0 Module
%
% Compile to .pdf with LaTeX (pdflatex)
% Requires Beamer package (latex-beamer)
%

\documentclass[xcolor=table]{beamer}
\usetheme{Warsaw}
\beamertemplatenavigationsymbolsempty
\setbeamertemplate{headline}{}
\useoutertheme{infolines}
\usepackage[english]{babel}
\usepackage[utf8]{inputenc}
\usepackage{graphics}
\usepackage{amssymb}
\usepackage{multirow}
\usepackage{xcolor}
\usepackage{framed}
\usepackage{hyperref}

\definecolor{shadecolor}{RGB}{180,180,180}

%%%%%%%%%%%%%%%%%%%%%%%%%%%%%%%%%%%%%%%%%%%%%%%%%%%%%%%%%%%%%%%%
% Python Highlighting
%%%%%%%%%%%%%%%%%%%%%%%%%%%%%%%%%%%%%%%%%%%%%%%%%%%%%%%%%%%%%%%%
\DeclareFixedFont{\ttb}{T1}{txtt}{bx}{n}{10}
\DeclareFixedFont{\ttm}{T1}{txtt}{m}{n}{10}

% Custom colors
\usepackage{color}
\definecolor{deepblue}{rgb}{0,0,0.5}
\definecolor{deepred}{rgb}{0.6,0,0}
\definecolor{deepgreen}{rgb}{0,0.5,0}
\definecolor{grey1}{rgb}{0.5,0.5,0.5}

\usepackage{listings}

% Python style for highlighting
\newcommand\pythonstyle{\lstset{
language=Python,
basicstyle=\ttm\small,
morekeywords={self},
keywordstyle=\ttb\color{deepblue},
emphstyle=\ttb\color{deepred},
stringstyle=\color{deepgreen},
commentstyle=\small\color{grey1}\ttm,
frame=single,
showstringspaces=false
}}

% Python environment
\lstnewenvironment{python}[1][]
{
\pythonstyle
\lstset{#1}
}
{}

% Python for inline
\newcommand\pythoninline[1]{{\pythonstyle\lstinline!#1!}}

\newcommand{\asignatura}{Session 9: Practice 2 -- Client0 Module}
\newcommand{\grado}{Programming in Network Environments}
\newcommand{\curso}{2025-2026}

\title[\asignatura]{\asignatura}
\subtitle{\grado}
\author{URJC}
\institute{GSyC}
\date{Course \curso}

\AtBeginSection[]
{
\begin{frame}<beamer>
\begin{center}
{\Huge \insertsection}
\end{center}
\end{frame}
}

\begin{document}

\frame{
\maketitle
}

\begin{frame}
\frametitle{Contents}
\tableofcontents
\end{frame}

%%%%%%%%%%%%%%%%%%%%%%%%%%%%%%%%%%%%%%%%%%%%%%%%%%%%%%%%%%%%%%%%
%%%%%%%%%%%%%%%%%%%%%%%%%%%%%%%%%%%%%%%%%%%%%%%%%%%%%%%%%%%%%%%%
% SECTION 1: Session 8 Review
%%%%%%%%%%%%%%%%%%%%%%%%%%%%%%%%%%%%%%%%%%%%%%%%%%%%%%%%%%%%%%%%
%%%%%%%%%%%%%%%%%%%%%%%%%%%%%%%%%%%%%%%%%%%%%%%%%%%%%%%%%%%%%%%%

\section{Review: Session 8 -- Client-Server 1}

\begin{frame}
\frametitle{What We Learned in Session 8 (1/2)}

\textbf{Session 8: Client-Server 1}:

\begin{itemize}
\item \textbf{Client-Server model}
  \begin{itemize}
  \item The \textbf{server} waits and attends requests from clients
  \item The \textbf{client} initiates the conversation
  \end{itemize}
\vspace{0.15cm}
\item \textbf{IP addresses}: numbers that identify network interfaces
  \begin{itemize}
  \item Find yours with \texttt{ifconfig} or \texttt{hostname -I}
  \item Test connectivity with \texttt{ping}
  \end{itemize}
\vspace{0.15cm}
\item \textbf{URLs and DNS}: human-readable names translated to IPs
\vspace{0.15cm}
\item \textbf{Ports}: identify which application to connect to
  \begin{itemize}
  \item Need both \textbf{(IP, port)} to reach a remote app
  \item Port 80 is reserved for web servers
  \end{itemize}
\end{itemize}

\end{frame}

\begin{frame}[fragile]
\frametitle{What We Learned in Session 8 (2/2)}

\textbf{Sockets}: the mechanism to communicate with remote programs.

\vspace{0.2cm}

\begin{enumerate}
\item \textbf{Create} the socket
\item \textbf{Connect} to (IP, Port)
\item \textbf{Send} data (as bytes, not strings)
\item \textbf{Receive} data from the server
\item \textbf{Close} the socket
\end{enumerate}

\vspace{0.2cm}

\begin{python}
import socket

s = socket.socket(socket.AF_INET,
                  socket.SOCK_STREAM)
s.connect((IP, PORT))
s.send(str.encode("Hello!"))
msg = s.recv(2048).decode("utf-8")
s.close()
\end{python}

\end{frame}

\begin{frame}[fragile]
\frametitle{Socket Lifecycle: Server vs Client}

\begingroup\small
\begin{columns}[T]

\begin{column}{0.45\textwidth}
\begin{center}
\textbf{SERVER}
\end{center}

\vspace{0.1cm}

\begin{enumerate}
\item \texttt{socket()} -- Create socket
\item \texttt{bind(IP, PORT)} -- Assign address
\item \texttt{listen()} -- Start listening
\item \texttt{accept()} -- Wait for client \\ {\scriptsize (blocks until connection)}
\item \texttt{recv()} -- Read client message
\item \texttt{send()} -- Send response
\item \texttt{close()} -- Close connection
\item Go to step 4 (loop)
\end{enumerate}
\end{column}

\begin{column}{0.45\textwidth}
\begin{center}
\textbf{CLIENT}
\end{center}

\vspace{0.1cm}

\begin{enumerate}
\item \texttt{socket()} -- Create socket
\item \texttt{connect(IP, PORT)} -- Connect to server
\item \texttt{send()} -- Send message
\item \texttt{recv()} -- Receive response
\item \texttt{close()} -- Close socket
\end{enumerate}

\vspace{0.3cm}

{\scriptsize The client does \textbf{not} need \texttt{bind()}, \texttt{listen()}, or \texttt{accept()} -- those are server-only operations.}
\end{column}

\end{columns}
\endgroup

\end{frame}

\begin{frame}
\frametitle{Session 8 Checklist}

\textbf{Make sure you have}:

\begin{itemize}
\item S08 folder in your repository with:
  \begin{itemize}
  \item \textbf{e1.txt} -- IP addresses table
  \item \textbf{e2.txt} -- Ping times table
  \item \textbf{e3.py} -- Chat client
  \item \textbf{server.py} -- Teacher's server running locally
  \end{itemize}
\item Everything pushed to your remote Github repo
\end{itemize}

\vspace{0.3cm}

\textbf{Key takeaways}:
\begin{itemize}
\item You can run \textbf{both client and server} on the same machine
\item Strings must be \textbf{encoded to bytes} before sending
\item Use \pythoninline{str.encode()} to send, \pythoninline{.decode("utf-8")} to receive
\end{itemize}

\end{frame}

%%%%%%%%%%%%%%%%%%%%%%%%%%%%%%%%%%%%%%%%%%%%%%%%%%%%%%%%%%%%%%%%
%%%%%%%%%%%%%%%%%%%%%%%%%%%%%%%%%%%%%%%%%%%%%%%%%%%%%%%%%%%%%%%%
% SECTION 2: Session 9 Overview
%%%%%%%%%%%%%%%%%%%%%%%%%%%%%%%%%%%%%%%%%%%%%%%%%%%%%%%%%%%%%%%%
%%%%%%%%%%%%%%%%%%%%%%%%%%%%%%%%%%%%%%%%%%%%%%%%%%%%%%%%%%%%%%%%

\section{Session 9: Practice 2 -- Client0 Module}

\begin{frame}
\frametitle{Session 9 Goals}

\textbf{In this session we will}:

\begin{itemize}
\item Create the \textbf{Client Class} in a module called \textbf{Client0.py}
\item Practice creating \textbf{clients} that talk to servers
\item Run \textbf{servers and clients} on the same machine
\item Get familiar with typical \textbf{socket and network errors}
\end{itemize}

\vspace{0.5cm}

\textbf{Time}: 2 hours

\vspace{0.3cm}

\textbf{Important}: Finish all S08 exercises before starting!

\end{frame}

\begin{frame}
\frametitle{The Client Class: Methods Overview}

\begin{center}
\begin{tabular}{|l|l|l|}
\hline
\textbf{Method} & \textbf{Parameters} & \textbf{Description} \\
\hline
\pythoninline{__init__(ip, port)} & IP: str, Port: int & Store IP and Port \\
\hline
\pythoninline{__str__()} & None & Print connection info \\
\hline
\pythoninline{ping()} & None & Print ``OK'' (test) \\
\hline
\pythoninline{talk(msg)} & msg: str & Send msg, return response \\
\hline
\end{tabular}
\end{center}

\vspace{0.3cm}

\textbf{Key idea}: Encapsulate all socket operations into a \textbf{reusable class}, so sending messages to a server becomes a single method call.

\vspace{0.3cm}

\pythoninline{__str__()} returns:\\
\texttt{"Connection to SERVER at \{ip\}, PORT: \{port\}"}

\end{frame}

\begin{frame}[fragile]
\frametitle{Exercises 1--3: Building the Client Class}

\textbf{Exercise 1}: \pythoninline{__init__()} and \pythoninline{ping()}
\begin{itemize}
\item Create \texttt{P02/Client0.py} (mark P02 as Sources Root)
\item \pythoninline{ping()} just prints ``OK!'' -- for testing the module
\end{itemize}

\vspace{0.2cm}

\textbf{Exercise 2}: \pythoninline{__str__()}
\begin{itemize}
\item Returns: \texttt{"Connection to SERVER at \{ip\}, PORT: \{port\}"}
\item No real connection yet -- just prints configuration
\end{itemize}

\vspace{0.2cm}

\textbf{Exercise 3}: \pythoninline{talk(msg)} -- the core method
\begin{itemize}
\item Creates socket, connects, sends, receives, closes
\item Returns the server's response as a string
\end{itemize}

\begin{python}
c = Client(IP, PORT)
response = c.talk("Testing!!!")
print(f"Response: {response}")
\end{python}

\end{frame}

\begin{frame}[fragile]
\frametitle{The talk() Method: Inside}

\pythoninline{talk(msg)} encapsulates the full socket workflow:

\vspace{0.2cm}

\begin{python}
def talk(self, msg):
    # Create the socket
    s = socket.socket(socket.AF_INET,
                      socket.SOCK_STREAM)
    # Connect to the server
    s.connect((self.ip, self.port))
    # Send data (encoded to bytes)
    s.send(str.encode(msg))
    # Receive response
    response = s.recv(2048).decode("utf-8")
    # Close the socket
    s.close()
    return response
\end{python}

\vspace{0.2cm}

\textbf{Remember}: Launch the server \textbf{first}, then run the client!

\end{frame}

\begin{frame}[fragile]
\frametitle{Exercises 4--5: Sending Genes to the Server}

\textbf{Exercise 4}: Send complete genes
\begin{itemize}
\item Use the \textbf{Seq Class} (from Practice 1) to load genes
\item Send U5, FRAT1, and ADA genes to the server (one at a time)
\item Use \pythoninline{str(s)} to get the sequence string from a Seq object
\end{itemize}

\vspace{0.3cm}

\textbf{Exercise 5}: Send gene \textbf{fragments}
\begin{itemize}
\item Take 5 fragments of 10 bases each from FRAT1
\item Send each fragment to the server using \pythoninline{talk()}
\item Print the fragments on the client console
\end{itemize}

\vspace{0.2cm}

\begin{python}
# Example: getting a fragment
gene_str = str(s)
fragment = gene_str[0:10]   # First 10 bases
\end{python}

\end{frame}

\begin{frame}
\frametitle{Exercise 6 and Deliverables}

\textbf{Exercise 6}: Two servers simultaneously
\begin{itemize}
\item Copy \texttt{server.py} as \texttt{server1.py} (port 8080) and \texttt{server2.py} (port 8081)
\item Take 10 fragments of 10 bases from FRAT1
\item \textbf{Odd} fragments (1,3,5,7,9) $\rightarrow$ server 1
\item \textbf{Even} fragments (2,4,6,8,10) $\rightarrow$ server 2
\end{itemize}

\vspace{0.4cm}

\textbf{P02 folder should contain}:

\begin{itemize}
\item \textbf{Client0.py} -- The module with the Client Class
\item \textbf{e1.py} through \textbf{e6.py} -- One program per exercise
\item \textbf{server1.py}, \textbf{server2.py} -- Server copies
\end{itemize}

\vspace{0.3cm}

\textbf{Remember}:
\begin{itemize}
\item Push everything to your remote Github repository
\item Make sure all exercises produce the expected output
\end{itemize}

\end{frame}

%%%%%%%%%%%%%%%%%%%%%%%%%%%%%%%%%%%%%%%%%%%%%%%%%%%%%%%%%%%%%%%%
%%%%%%%%%%%%%%%%%%%%%%%%%%%%%%%%%%%%%%%%%%%%%%%%%%%%%%%%%%%%%%%%
% SECTION 3: End
%%%%%%%%%%%%%%%%%%%%%%%%%%%%%%%%%%%%%%%%%%%%%%%%%%%%%%%%%%%%%%%%
%%%%%%%%%%%%%%%%%%%%%%%%%%%%%%%%%%%%%%%%%%%%%%%%%%%%%%%%%%%%%%%%

\section{Let's start!}

\begin{frame}
\begin{center}
\Huge{Questions?}

\vspace{1cm}

\Large{Time to code!}
\end{center}
\end{frame}

\end{document}
